% 第22章 重积分：题目解答
% 仅解答原章中的思考题、练习题、参考题；原书已有解答的普通例题不重复。

\chapter*{第22章\quad 重积分习题解答}
\addcontentsline{toc}{chapter}{第22章\quad 重积分习题解答}
\subsection*{22.1.3 思考题}

\paragraph{第1题}
\begin{xsolution}
若 $f,g$ 在 $D$ 上可积，则二者均有界。由可积判别准则，$f,g$ 的不连续点集分别为零测度集 $E_f,E_g$。在 $D\setminus(E_f\cup E_g)$ 上，$f,g$ 都连续，因此 $fg$ 连续。故 $fg$ 的不连续点集包含于 $E_f\cup E_g$，仍为零测度集，所以 $fg$ 在 $D$ 上可积。

题目中关于倒数的断言按原条件并不成立。取 $D=[0,1]\times[0,1]$，定义
\[
  f(x,y)=
  \begin{cases}
    x,&x>0,\\
    1,&x=0.
  \end{cases}
\]
$f$ 有界，且仅在直线 $x=0$ 上不连续，所以在 $D$ 上可积；又处处有 $f\ne0$。但是
\[
  \frac1{f(x,y)}=
  \begin{cases}
    1/x,&x>0,\\
    1,&x=0
  \end{cases}
\]
在 $D$ 上无界，因而不是通常意义下的 Riemann 可积函数。

若补充条件 $\abs{f(x,y)}\ge m>0$，则 $1/f$ 有界，而且在 $f$ 的每个连续点处连续，故其不连续点集仍包含于 $f$ 的不连续点集，从而 $1/f$ 可积。
\end{xsolution}

\paragraph{第2题}
\begin{xsolution}
设 $u(D)\subset[m,M]$。因为 $u$ 可积，所以 $u$ 有界，且其不连续点集 $E$ 为零测度集。若 $f$ 在 $[m,M]$ 上连续，则在 $D\setminus E$ 的每一点 $(x_0,y_0)$，$u$ 连续，复合函数连续性给出
\[
  f(u(x,y))\longrightarrow f(u(x_0,y_0)).
\]
因此 $f\circ u$ 的不连续点集包含于 $E$。又 $f$ 在紧区间上有界，所以 $f(u(x,y))$ 有界。由可积判别准则，$f(u(x,y))$ 在 $D$ 上可积。

若 $f$ 仅仅可积，结论不成立。取 $D=[0,1]^2$，令 Thomae 函数
\[
  T(x)=
  \begin{cases}
    1/q,&x=p/q\text{ 为既约有理数},\\
    0,&x\text{ 为无理数},
  \end{cases}
\]
并令 $u(x,y)=T(x)$。$u$ 的不连续点集为 $\QQ\cap[0,1]$ 与 $[0,1]$ 的直积，是零测度集，所以 $u$ 可积。再定义
\[
  f(t)=
  \begin{cases}
    1,&t=0,\\
    0,&t\ne0.
  \end{cases}
\]
$f$ 在 $[0,1]$ 上只有 $t=0$ 一个不连续点，故可积。然而
\[
  f(u(x,y))=
  \begin{cases}
    1,&x\text{ 为无理数},\\
    0,&x\text{ 为有理数},
  \end{cases}
\]
在 $D$ 的每一点都不连续，因而不可积。
\end{xsolution}

\paragraph{第3题}
\begin{xsolution}
先设 $f=g$ 除了一个零面积集 $E$ 外处处成立，并令 $h=f-g$。由于 $f,g$ 有界，存在 $M>0$ 使 $\abs{h}\le M$。对任意 $\varepsilon>0$，可用有限个矩形覆盖 $E$，使这些矩形的总面积小于 $\varepsilon$。在覆盖之外 $h=0$，所以对应任一足够细的分划，$h$ 的上下和之差只来自这些矩形，且不超过 $2M\varepsilon$。由 $\varepsilon$ 的任意性，$h$ 可积且
\[
  \iint_D h(x,y)\,\dd x\,\dd y=0.
\]
于是 $f$ 与 $g$ 有相同的可积性；可积时
\[
  \iint_Df\,\dd x\,\dd y=\iint_Dg\,\dd x\,\dd y.
\]

若把“零面积集”换成一般的“零测度集”，结论不再成立。取 $D=[0,1]^2$，令 $f\equiv0$，而
\[
  g(x,y)=\mathbf 1_{(\QQ\cap[0,1])^2}(x,y).
\]
两函数仅在可列集 $(\QQ\cap[0,1])^2$ 上不同，这个集合是零测度集；但 $g$ 在每一点附近同时取得 $0$ 和 $1$，故处处不连续，不可积，而 $f$ 可积。
\end{xsolution}

\paragraph{第4题}
\begin{xsolution}
\begin{enumerate}[label=(\arabic*)]
  \item 函数
  \[
    f(x,y)=\sin\frac1{(x^2-1)^2+(y^2-1)^2}
  \]
  在正方形 $[-1,1]^2$ 内有界，只有四个角点 $(\pm1,\pm1)$ 是奇点。在这些点任意补定义函数值后，所得函数的不连续点集至多为这四点，因而是零测度集。故函数在 $[-1,1]^2$ 上可积；对原函数按题目所述约定，其积分就是去掉四个角点后的积分。

  \item 函数
  \[
    f(x,y)=\arctan\frac1{y-x^2}
  \]
  在 $[0,1]^2$ 上有界。除抛物线弧 $y=x^2$ 外处处连续，而该曲线面积为零。在曲线上任意补定义后，不连续点集包含于这条零面积曲线，故函数可积。于是按题目约定，原函数在 $[0,1]^2$ 上可积。
\end{enumerate}
\end{xsolution}

\subsection*{22.1.4 练习题}

\paragraph{第1题}
\begin{xsolution}
由
\[
  \max\{f,g\}=\frac12\bigl(f+g+\abs{f-g}\bigr)
\]
即可。$f-g$ 可积；若 $u$ 可积，则 $\abs{u}$ 的不连续点只能出现在 $u$ 的不连续点处，因此 $\abs{u}$ 也可积。于是 $\abs{f-g}$ 可积，从而 $h=\max\{f,g\}$ 可积。
\end{xsolution}

\paragraph{第2题}
\begin{xsolution}
记 $B_\rho(p_0)=\{(x,y)\mid(x-x_0)^2+(y-y_0)^2\le\rho^2\}$。由 $f$ 在 $p_0$ 连续，对任意 $\varepsilon>0$，存在 $\delta>0$，当 $0<\rho<\delta$ 且 $(x,y)\in B_\rho(p_0)$ 时有
\[
  \abs{f(x,y)-f(x_0,y_0)}<\varepsilon.
\]
因此
\begin{align*}
 &\abs*{\frac1{\pi\rho^2}\iint_{B_\rho(p_0)}f(x,y)\,\dd x\,\dd y-f(x_0,y_0)}\\
 &\qquad\le\frac1{\pi\rho^2}\iint_{B_\rho(p_0)}\abs{f(x,y)-f(x_0,y_0)}\,\dd x\,\dd y<\varepsilon.
\end{align*}
故极限为
\[
  \boxed{f(x_0,y_0)}.
\]
\end{xsolution}

\paragraph{第3题}
\begin{xsolution}
作线性变换
\[
  u=x+y,\qquad v=x-y,
\]
则
\[
  x=\frac{u+v}{2},\qquad y=\frac{u-v}{2},\qquad
  \abs*{\frac{\partial(x,y)}{\partial(u,v)}}=\frac12.
\]
并且
\[
  \abs{x}+\abs{y}=\max\{\abs{u},\abs{v}\},
\]
故菱形 $\abs{x}+\abs{y}\le1$ 被变成正方形 $-1\le u,v\le1$。于是
\begin{align*}
  \iint_{\abs{x}+\abs{y}\le1}f(x+y)\,\dd x\,\dd y
  &=\frac12\int_{-1}^1\int_{-1}^1 f(u)\,\dd v\,\dd u\\
  &=\int_{-1}^1f(u)\,\dd u.
\end{align*}
\end{xsolution}

\paragraph{第4题}
\begin{xsolution}
在第一象限作变换
\[
  u=xy,\qquad v=\frac yx.
\]
区域的四条边分别变为 $u=1,u=2,v=1,v=4$。由
\[
  x=\sqrt{\frac uv},\qquad y=\sqrt{uv},
\]
得
\[
  \abs*{\frac{\partial(x,y)}{\partial(u,v)}}=\frac1{2v}.
\]
故
\begin{align*}
  \iint_Df(xy)\,\dd x\,\dd y
  &=\int_1^2\int_1^4 f(u)\frac1{2v}\,\dd v\,\dd u\\
  &=\frac12\ln4\int_1^2f(u)\,\dd u
  =\ln2\int_1^2f(u)\,\dd u.
\end{align*}
\end{xsolution}

\subsection*{22.2.4 练习题}

\paragraph{第1题}
\begin{xsolution}
原积分区域位于第一象限，由直线 $y=Rx$、圆 $x^2+y^2=R^2$ 和 $x$ 轴围成。直线与圆的交点为
\[
  \left(\frac{R}{\sqrt{1+R^2}},\frac{R^2}{\sqrt{1+R^2}}\right).
\]
固定 $y$ 后，有
\[
  \frac yR\le x\le\sqrt{R^2-y^2},
  \qquad0\le y\le\frac{R^2}{\sqrt{1+R^2}}.
\]
所以
\[
  \boxed{
  I=\int_0^{R^2/\sqrt{1+R^2}}\dd y
  \int_{y/R}^{\sqrt{R^2-y^2}}f(x,y)\,\dd x }.
\]
\end{xsolution}

\paragraph{第2题}
\begin{xsolution}
积分区域为 $0\le x\le t\le1$。交换积分次序得
\begin{align*}
  \int_0^1\dd x\int_x^1f(t)\,\dd t
  &=\int_0^1\dd t\int_0^t f(t)\,\dd x\\
  &=\int_0^1t f(t)\,\dd t.
\end{align*}
\end{xsolution}

\paragraph{第3题}
\begin{xsolution}
三条边围成的三角形可写成
\[
  0\le x\le\pi,\qquad \pi-x\le y\le\pi.
\]
故
\begin{align*}
  \iint_D\frac{\sin x}{x}\,\dd x\,\dd y
  &=\int_0^\pi\frac{\sin x}{x}\bigl[\pi-(\pi-x)\bigr]\,\dd x\\
  &=\int_0^\pi\sin x\,\dd x=\boxed{2}.
\end{align*}
$x=0$ 处按连续延拓理解即可。
\end{xsolution}

\paragraph{第4题}
\begin{xsolution}
曲线 $y=x^2$ 与直线 $x+y=2$ 在第一象限交于 $(1,1)$。因此
\[
  D=\{0\le x\le1,0\le y\le x^2\}
  \cup\{1\le x\le2,0\le y\le2-x\}.
\]
于是
\begin{align*}
 I&=\int_0^1\int_0^{x^2}(x^2+y^2)\,\dd y\,\dd x
 +\int_1^2\int_0^{2-x}(x^2+y^2)\,\dd y\,\dd x\\
 &=\int_0^1\left(x^4+\frac{x^6}{3}\right)\dd x
 +\int_1^2\left[x^2(2-x)+\frac{(2-x)^3}{3}\right]\dd x\\
 &=\frac{26}{105}+1=\boxed{\frac{131}{105}}.
\end{align*}
\end{xsolution}

\paragraph{第5题}
\begin{xsolution}
在极坐标下，边界曲线变为
\[
  r^4=a^2r^2\cos2\theta,
  \qquad r^2=a^2\cos2\theta.
\]
它有关于 $x$ 轴正、负方向的两个全等叶瓣；在叶瓣内 $x^2-y^2=r^2\cos2\theta\ge0$，故立体高度为 $r^2\cos2\theta$。于是
\begin{align*}
 V&=2\int_{-\pi/4}^{\pi/4}\int_0^{a\sqrt{\cos2\theta}}
 r^2\cos2\theta\,r\,\dd r\,\dd\theta\\
 &=\frac{a^4}{2}\int_{-\pi/4}^{\pi/4}\cos^3 2\theta\,\dd\theta
 =\frac{a^4}{4}\int_{-\pi/2}^{\pi/2}\cos^3u\,\dd u\\
 &=\boxed{\frac{a^4}{3}}.
\end{align*}
\end{xsolution}

\paragraph{第6题}
\begin{xsolution}
区域位于第一象限。作变换
\[
  u=\frac{x^2}{y},\qquad v=\frac{y^2}{x}.
\]
四族边界分别为 $u=a,u=b,v=p,v=q$，而
\[
  uv=xy,
  \qquad
  \frac{\partial(u,v)}{\partial(x,y)}=3.
\]
故 $\dd x\,\dd y=\frac13\dd u\,\dd v$，且
\begin{align*}
 I&=\frac13\int_a^b\int_p^q u\sin(uv)\,\dd v\,\dd u\\
 &=\frac13\int_a^b\bigl[\cos(pu)-\cos(qu)\bigr]\,\dd u\\
 &=\boxed{\frac13\left[
 \frac{\sin(bp)-\sin(ap)}p-
 \frac{\sin(bq)-\sin(aq)}q\right]}.
\end{align*}
\end{xsolution}

\paragraph{第7题}
\begin{xsolution}
原积分区域为 $a\le x\le y\le b$。交换积分次序：$a\le x\le b$，$x\le y\le b$。于是（$n$ 为非负整数）
\begin{align*}
 \int_a^b\dd y\int_a^y(y-x)^nf(x)\,\dd x
 &=\int_a^b f(x)\,\dd x\int_x^b(y-x)^n\,\dd y\\
 &=\frac1{n+1}\int_a^b(b-x)^{n+1}f(x)\,\dd x.
\end{align*}
\end{xsolution}

\paragraph{第8题}
\begin{xsolution}
在第一象限，两圆分别为
\[
  r=a,
  \qquad r=2a\cos\theta.
\]
它们在 $\theta=\pi/3$ 相交。由 $y$ 轴封闭的区域为
\[
  \frac\pi3\le\theta\le\frac\pi2,
  \qquad2a\cos\theta\le r\le a.
\]
被积函数为 $r$，故
\begin{align*}
 I&=\int_{\pi/3}^{\pi/2}\int_{2a\cos\theta}^{a}r^2\,\dd r\,\dd\theta\\
 &=\frac{a^3}{3}\left[\frac\pi6-8\int_{\pi/3}^{\pi/2}\cos^3\theta\,\dd\theta\right]\\
 &=\boxed{\frac{a^3}{18}\left(\pi+18\sqrt3-32\right)}.
\end{align*}
\end{xsolution}

\paragraph{第9题}
\begin{xsolution}
令
\[
  u=x+y,\qquad v=\frac yx.
\]
则 $p\le u\le q$，$a\le v\le b$，且
\[
  x=\frac{u}{1+v},\qquad y=\frac{uv}{1+v},
  \qquad
  \abs*{\frac{\partial(x,y)}{\partial(u,v)}}=\frac{u}{(1+v)^2}.
\]
所以所求面积
\begin{align*}
 S&=\int_p^q\int_a^b\frac{u}{(1+v)^2}\,\dd v\,\dd u\\
 &=\frac{q^2-p^2}{2}\left(\frac1{1+a}-\frac1{1+b}\right)\\
 &=\boxed{\frac{(q^2-p^2)(b-a)}{2(1+a)(1+b)}}.
\end{align*}
\end{xsolution}

\paragraph{第10题}
\begin{xsolution}
令
\[
  u=\sqrt[4]{x/a},\qquad v=\sqrt[4]{y/b}.
\]
则 $x=au^4,y=bv^4$，区域变为
\[
  u\ge0,\quad v\ge0,\quad u+v\le1,
\]
且
\[
  \abs*{\frac{\partial(x,y)}{\partial(u,v)}}=16ab\,u^3v^3.
\]
故
\begin{align*}
 S&=16ab\int_0^1u^3\,\dd u\int_0^{1-u}v^3\,\dd v\\
 &=4ab\int_0^1u^3(1-u)^4\,\dd u
 =4ab\cdot\frac1{280}=\boxed{\frac{ab}{70}}.
\end{align*}
\end{xsolution}

\paragraph{第11题}
\begin{xsolution}
区域由关于正、负 $x$ 轴的两个全等扇形状部分组成。极坐标下，角度范围为
\[
 -\frac\pi4\le\theta\le\frac\pi4,
 \qquad
 \frac{3\pi}{4}\le\theta\le\frac{5\pi}{4},
\]
且 $r\abs{\cos\theta}\le1$。对 $0\le r\le1$，可取的角度总长为 $\pi$；对 $1\le r\le\sqrt2$，可取的角度总长为
\[
  \pi-4\arccos\frac1r.
\]
因此
\[
 \boxed{
 \iint_Df(\sqrt{x^2+y^2})\,\dd x\,\dd y
 =\pi\int_0^{\sqrt2}f(r)r\,\dd r
 -4\int_1^{\sqrt2}\arccos\frac1r\,f(r)r\,\dd r }.
\]
\end{xsolution}

\paragraph{第12题}
\begin{xsolution}
由 $xy=1$ 与 $x^2+y^2=4$ 围成的整个区域关于原点对称，而被积函数 $x$ 满足
\[
  x(-x,-y)=-x(x,y).
\]
故两对称部分的积分互相抵消，
\[
  \boxed{\iint_Dx\,\dd x\,\dd y=0}.
\]
\end{xsolution}

\paragraph{第13题}
\begin{xsolution}
记
\[
 a=x_u=y_v,
 \qquad b=x_v=-y_u.
\]
正则性给出
\[
 J=\frac{\partial(x,y)}{\partial(u,v)}=x_uy_v-x_vy_u=a^2+b^2>0.
\]
由链式法则
\[
 f_u=f_xa-f_yb,
 \qquad
 f_v=f_xb+f_ya,
\]
故
\[
 f_u^2+f_v^2=(a^2+b^2)(f_x^2+f_y^2)=J(f_x^2+f_y^2).
\]
又 $\dd x\,\dd y=J\,\dd u\,\dd v$，所以
\[
 \iint_D(f_x^2+f_y^2)\,\dd x\,\dd y
 =\iint_\Omega(f_u^2+f_v^2)\,\dd u\,\dd v.
\]
\end{xsolution}

\paragraph{第14题}
\begin{xsolution}
积分区域满足 $x\ge y\ge0$ 且 $x^2+y^2\le4$，即极坐标下
\[
 0\le\theta\le\frac\pi4,
 \qquad0\le r\le2.
\]
于是
\begin{align*}
 I&=\int_0^{\pi/4}\int_0^2\frac{r}{\sqrt{1+r^2}}\,\dd r\,\dd\theta\\
 &=\frac\pi4\left[\sqrt{1+r^2}\right]_0^2
 =\boxed{\frac\pi4(\sqrt5-1)}.
\end{align*}
\end{xsolution}

\paragraph{第15题}
\begin{xsolution}
由四象限对称性，
\begin{align*}
 I&=4\int_0^1\int_0^{1-x}(\sqrt{xy}+xy)\,\dd y\,\dd x\\
 &=\frac83\int_0^1x^{1/2}(1-x)^{3/2}\,\dd x
 +2\int_0^1x(1-x)^2\,\dd x\\
 &=\frac83B\left(\frac32,\frac52\right)+\frac16
 =\frac\pi6+\frac16.
\end{align*}
因此
\[
 I=\frac{\pi+1}{6}<\frac32,
\]
所要求的不等式成立。
\end{xsolution}

\paragraph{第16题}
\begin{xsolution}
固定 $y$，令 $a=y^2\in[0,1]$，则
\[
 \int_0^1\abs{x-y^2}\,\dd x
 =\int_0^a(a-x)\,\dd x+\int_a^1(x-a)\,\dd x
 =a^2-a+\frac12.
\]
故
\begin{align*}
 I&=\int_{-1}^1\left(y^4-y^2+\frac12\right)\dd y\\
 &=2\left(\frac15-\frac13+\frac12\right)
 =\boxed{\frac{11}{15}}.
\end{align*}
\end{xsolution}

\paragraph{第17题}
\begin{xsolution}
三角形可写成
\[
 1\le x\le2,
 \qquad1\le y\le x.
\]
在其中 $\ln(x/y^2)$ 的变号曲线为 $y=\sqrt x$。于是
\begin{align*}
 I&=\int_1^2\left[
 \int_1^{\sqrt x}\ln\frac{x}{y^2}\,\dd y
 -\int_{\sqrt x}^{x}\ln\frac{x}{y^2}\,\dd y
 \right]\dd x\\
 &=\boxed{\frac{64\sqrt2-89}{12}}.
\end{align*}
例如使用原函数
\[
 \int\ln\frac{x}{y^2}\,\dd y
 =y\ln x-2y\ln y+2y
\]
代入三个端点即可得到上述结果。
\end{xsolution}

\subsection*{22.3.5 练习题}

\paragraph{第1题}
\begin{xsolution}
原积分区域为
\[
 x\ge0,\quad z\ge0,\quad y\ge0,\quad x+z+y\le1.
\]
在最内层令 $t=1-y-z$，则 $y=1-z-t$，且 $1-y=z+t$。再交换 $x,z,t$ 的次序，可写成
\[
 0\le t\le1,
 \qquad0\le x\le t,
 \qquad0\le z\le1-t.
\]
于是
\begin{align*}
 I&=\int_0^1\ee^{-t^2}\,\dd t\int_0^t\dd x\int_0^{1-t}(z+t)\,\dd z\\
 &=\frac12\int_0^1t(1-t^2)\ee^{-t^2}\,\dd t.
\end{align*}
令 $u=t^2$，得
\[
 I=\frac14\int_0^1(1-u)\ee^{-u}\,\dd u
 =\boxed{\frac1{4\ee}}.
\]
\end{xsolution}

\paragraph{第2题}
\begin{xsolution}
在 $xOy$ 平面上，
\[
 (x-1)^2+y^2\le1,
 \qquad y\le0.
\]
作柱坐标变换
\[
 x=r\cos\theta,
 \qquad y=r\sin\theta,
\]
则该半圆区域为
\[
 -\frac\pi2\le\theta\le0,
 \qquad0\le r\le2\cos\theta.
\]
又 $0\le z\le x=r\cos\theta$。所以原累次积分化为
\[
 \boxed{
 \int_{-\pi/2}^{0}\dd\theta
 \int_0^{2\cos\theta}r\,\dd r
 \int_0^{r\cos\theta}
 f(r\cos\theta,r\sin\theta,z)\,\dd z }.
\]
\end{xsolution}

\paragraph{第3题}
\begin{xsolution}
旋转曲面为抛物面
\[
 x^2+y^2=2z.
\]
在柱坐标下，区域为
\[
 2\le z\le8,
 \qquad0\le r\le\sqrt{2z},
 \qquad0\le\theta\le2\pi.
\]
因此
\begin{align*}
 I&=\int_2^8\int_0^{2\pi}\int_0^{\sqrt{2z}}r^2\,r\,\dd r\,\dd\theta\,\dd z\\
 &=2\pi\int_2^8 z^2\,\dd z
 =\boxed{336\pi}.
\end{align*}
\end{xsolution}

\paragraph{第4题}
\begin{xsolution}
用柱坐标 $x=r\cos\theta,y=r\sin\theta$。由 $x,y\ge0$ 有 $0\le\theta\le\pi/2$。记
\[
 s=\sqrt{4-r^2}.
\]
两个球的公共部分在给定 $r$ 处满足
\[
 2-s\le z\le s,
\]
故必须 $s\ge1$，即 $0\le r\le\sqrt3$。于是
\begin{align*}
 I&=\int_0^{\pi/2}\cos\theta\sin\theta\,\dd\theta
 \int_0^{\sqrt3}r^3\,\dd r\int_{2-s}^{s}z\,\dd z\\
 &=\int_0^{\sqrt3}r^3(\sqrt{4-r^2}-1)\,\dd r
 =\boxed{\frac{53}{60}}.
\end{align*}
\end{xsolution}

\paragraph{第5题}
\begin{xsolution}
设球心为 $O$，球外固定点为 $P$，$A=OP>R$。以 $OP$ 为极轴，在以 $O$ 为原点的球坐标中，球内一点到 $P$ 的距离为
\[
 r=\sqrt{A^2+s^2-2As\cos\varphi},
\]
其中 $0\le s\le R$。于是
\begin{align*}
 I&=\int_0^R s^2\,\dd s\int_0^{2\pi}\dd\theta
 \int_0^\pi\frac{\sin\varphi\,\dd\varphi}
 {\sqrt{A^2+s^2-2As\cos\varphi}}.
\end{align*}
因 $A>s$，令 $u=\cos\varphi$ 可得
\[
 \int_0^\pi\frac{\sin\varphi\,\dd\varphi}
 {\sqrt{A^2+s^2-2As\cos\varphi}}=\frac2A.
\]
故
\[
 I=\frac{4\pi}{A}\int_0^Rs^2\,\dd s
 =\boxed{\frac{4\pi R^3}{3A}}.
\]
\end{xsolution}

\paragraph{第6题}
\begin{xsolution}
采用球坐标
\[
 x=\rho\sin\varphi\cos\theta,
 \quad y=\rho\sin\varphi\sin\theta,
 \quad z=\rho\cos\varphi.
\]
由 $z\ge0$ 有 $0\le\varphi\le\pi/2$。曲面方程化为
\[
 \rho^2=\frac{a^2}{2}\sin^2\varphi\sin2\theta.
\]
故 $\sin2\theta\ge0$，即
\[
 \theta\in\left[0,\frac\pi2\right]\cup
 \left[\pi,\frac{3\pi}{2}\right],
\]
且
\[
 0\le\rho\le\frac a{\sqrt2}\sin\varphi\sqrt{\sin2\theta}.
\]
同时
\[
 \frac{xyz}{x^2+y^2}=\rho\cos\varphi\cos\theta\sin\theta.
\]
于是
\begin{align*}
 I&=\frac{a^4}{32}
 \int_0^{\pi/2}\sin^5\varphi\cos\varphi\,\dd\varphi
 \left(
 \int_0^{\pi/2}\sin^32\theta\,\dd\theta+
 \int_\pi^{3\pi/2}\sin^32\theta\,\dd\theta
 \right)\\
 &=\frac{a^4}{32}\cdot\frac16\cdot\frac43
 =\boxed{\frac{a^4}{144}}.
\end{align*}
\end{xsolution}

\paragraph{第7题}
\begin{xsolution}
被积函数只依赖四维半径 $r=\sqrt{x^2+y^2+z^2+u^2}$。三维单位球面 $S^3$ 的面积为 $2\pi^2$，第一卦限占 $1/16$，故角向因子为 $\pi^2/8$。于是
\[
 I=\frac{\pi^2}{8}\int_0^1r^3
 \sqrt{\frac{1-r^2}{1+r^2}}\,\dd r.
\]
令 $t=r^2$，再令 $s=\sqrt{(1-t)/(1+t)}$，则
\[
 \int_0^1t\sqrt{\frac{1-t}{1+t}}\,\dd t
 =4\int_0^1\frac{s^2(1-s^2)}{(1+s^2)^3}\,\dd s
 =1-\frac\pi4.
\]
因此
\[
 \boxed{I=\frac{\pi^2}{16}\left(1-\frac\pi4\right)
 =\frac{\pi^2(4-\pi)}{64}}.
\]
\end{xsolution}

\paragraph{第8题}
\begin{xsolution}
球坐标下
\[
 F(t)=4\pi\int_0^t f(r^2)r^2\,\dd r.
\]
由微积分基本定理
\[
 F'(t)=4\pi t^2f(t^2).
\]
已知 $f(1)=1$，所以
\[
 \boxed{F'(1)=4\pi}.
\]
\end{xsolution}

\paragraph{第9题}
\begin{xsolution}
在球坐标中，条件 $z\ge\sqrt{x^2+y^2}$ 等价于
\[
 0\le\varphi\le\frac\pi4,
\]
而球壳条件为 $2\le r\le4$。被积函数 $f=r$。由于角向因子在分子与分母中相同，积分平均值等于
\[
 \frac{\int_2^4r^3\,\dd r}{\int_2^4r^2\,\dd r}
 =\frac{(4^4-2^4)/4}{(4^3-2^3)/3}
 =\boxed{\frac{45}{14}}.
\]
\end{xsolution}

\paragraph{第10题}
\begin{xsolution}
由 $xy=1,xy=2,y=3x,y=4x$ 围出的集合有第一、第三象限两个全等的连通分支。题目所说的“区域”按通常约定取其中一个连通分支；取第一象限分支。先对 $z$ 积分：
\[
 I=\frac12\iint_Dx^2y^2(x^2+y^2)^2\,\dd x\,\dd y.
\]
令
\[
 u=xy,
 \qquad v=\frac yx.
\]
则 $1\le u\le2,3\le v\le4$，并有
\[
 x^2+y^2=u\left(v+\frac1v\right),
 \qquad
 \dd x\,\dd y=\frac1{2v}\,\dd u\,\dd v.
\]
因此
\begin{align*}
 I&=\frac14\int_1^2u^4\,\dd u
 \int_3^4\left(v+\frac2v+\frac1{v^3}\right)\dd v\\
 &=\frac14\cdot\frac{31}{5}
 \left(\frac72+2\ln\frac43+\frac7{288}\right)\\
 &=\boxed{\frac{6293}{1152}+\frac{31}{10}\ln\frac43}.
\end{align*}
若把两个对称连通分支的并集都计入，则积分值为上述结果的 $2$ 倍。
\end{xsolution}

\paragraph{第11题}
\begin{xsolution}
令
\[
 k=\sqrt{a^2+b^2+c^2}>0.
\]
取正交变换把向量 $(a,b,c)$ 变到新坐标轴的第三轴方向。单位球在正交变换下保持不变，且
\[
 ax+by+cz=kw,
\]
其中 $w$ 为新坐标。固定 $w\in[-1,1]$ 时，截面是半径 $\sqrt{1-w^2}$ 的圆盘，故
\begin{align*}
 I&=\pi\int_{-1}^1(1-w^2)\cos(kw)\,\dd w\\
 &=\boxed{\frac{4\pi}{k^3}(\sin k-k\cos k)}.
\end{align*}
\end{xsolution}

\subsection*{22.4.4 练习题}

\paragraph{第1题}
\begin{xsolution}
以下按通常情形 $p,q>0$ 讨论。
\begin{enumerate}[label=(\arabic*)]
  \item 积分可分离为
  \[
   \left(\int_{-\infty}^{+\infty}\frac{\dd x}{1+\abs{x}^p}\right)
   \left(\int_{-\infty}^{+\infty}\frac{\dd y}{1+\abs{y}^q}\right).
  \]
  一维积分分别在且仅在 $p>1,q>1$ 时收敛。因此原广义积分
  \[
    \boxed{\text{当且仅当 }p>1,\ q>1\text{ 时收敛}}.
  \]

  \item 只需考察无穷远处。第一象限中，对 $x$ 足够大，
  \[
   \int_0^{+\infty}\frac{\dd y}{x^p+y^q}
   =x^{p/q-p}\int_0^{+\infty}\frac{\dd t}{1+t^q}.
  \]
  右端的 $t$ 积分要求 $q>1$；随后关于 $x$ 的尾积分收敛要求
  \[
    p-\frac pq>1
    \quad\Longleftrightarrow\quad
    \frac1p+\frac1q<1.
  \]
  这个条件本身已推出 $p,q>1$，并且四个象限完全相同。故
  \[
    \boxed{\text{当且仅当 }\frac1p+\frac1q<1\text{ 时收敛}}.
  \]

  \item 作变换
  \[
    u=x+y,
    \qquad v=x-y.
  \]
  在固定条带 $1\le u\le2$ 内，$v$ 可取任意实数，而
  \[
   \sin x\sin y
   =\frac12\left(\cos v-\cos u\right).
  \]
  因而在该条带上 $u^{-p}$ 有上下正界，而 $\abs{\cos v-\cos u}$ 关于 $v$ 的积分为无穷大。于是绝对积分发散。根据本章所述广义重积分“收敛当且仅当绝对收敛”的性质，原积分对任意 $p$ 都发散：
  \[
    \boxed{\text{对所有 }p\text{ 均发散}}.
  \]
\end{enumerate}
\end{xsolution}

\paragraph{第2题}
\begin{xsolution}
令
\[
 L=\lim_{n\to\infty}\iint_{D_n}f,
\]
允许 $L=+\infty$。由于 $f\ge0$ 且 $D_n$ 单调增加，序列中的积分单调增加，所以此极限总在 $[0,+\infty]$ 中存在。

对 $D$ 的任意有界闭子区域 $E$，令 $E_n=E\cap D_n$。则 $E_n$ 单调增加且并集为 $E$。由提示中的有界区域结论，
\[
 \iint_Ef=\lim_{n\to\infty}\iint_{E_n}f
 \le\lim_{n\to\infty}\iint_{D_n}f=L.
\]
因此所有有界子区域上的积分上确界不超过 $L$。反过来，每个 $D_n$ 本身就是 $D$ 的有界子区域，故该上确界不小于 $\iint_{D_n}f$，对 $n\to\infty$ 得不小于 $L$。所以二者相等。

因此 $L<+\infty$ 时 $f$ 在 $D$ 上广义可积且积分值为 $L$；$L=+\infty$ 时两边同时发散。这就证明了题中等式及“同时有意义或同时无意义”的结论。
\end{xsolution}

\paragraph{第3题}
\begin{xsolution}
作变换
\[
 x=t\sqrt y,
 \qquad y=y.
\]
由 $y\ge x^2+1$ 得
\[
 -1<t<1,
 \qquad y\ge\frac1{1-t^2},
\]
且 $\dd x\,\dd y=\sqrt y\,\dd t\,\dd y$。于是
\begin{align*}
 I&=\int_{-1}^1\frac{\dd t}{1+t^4}
 \int_{1/(1-t^2)}^{+\infty}y^{-3/2}\,\dd y\\
 &=4\int_0^1\frac{\sqrt{1-t^2}}{1+t^4}\,\dd t.
\end{align*}
令 $t=\sin\theta$，再令 $u=\tan\theta$，得
\[
 I=4\int_0^{+\infty}\frac{\dd u}{1+2u^2+2u^4}.
\]
令 $u=2^{-1/4}s$，利用
\[
 \int_0^{+\infty}\frac{\dd s}{s^4+\alpha s^2+1}
 =\frac\pi{2\sqrt{2+\alpha}},
 \qquad \alpha>-2,
\]
可得
\[
 \boxed{I=\frac{2^{3/4}\pi}{\sqrt{2+\sqrt2}}}.
\]
\end{xsolution}

\paragraph{第4题}
\begin{xsolution}
以下按 $p>0$ 讨论。
\begin{enumerate}[label=(\arabic*)]
  \item 唯一需要检查的是原点附近。由 $\abs{y}\le x^2$，当 $x\ne0$ 时
  \[
   \int_{-x^2}^{x^2}\frac{\dd y}{x^2+y^2}
   \le\int_{-x^2}^{x^2}\frac{\dd y}{x^2}=2.
  \]
  再对 $x$ 在原点附近积分即有限，故
  \[
    \boxed{\text{收敛}}.
  \]

  \item 二次型
  \[
    x^2+xy+y^2
  \]
  正定，所以存在常数 $c_1,c_2>0$ 使
  \[
    c_1(x^2+y^2)\le x^2+xy+y^2\le c_2(x^2+y^2).
  \]
  因此原点附近的积分与
  \[
    \int_0^1r^{1-2p}\,\dd r
  \]
  同敛散，故
  \[
    \boxed{0<p<1\text{ 时收敛， }p\ge1\text{ 时发散}}.
  \]

  \item 极坐标下
  \[
   I=2\pi\int_0^1\frac{r\,\dd r}{(1-r^2)^p}
   =\pi\int_0^1(1-u)^{-p}\,\dd u,
  \]
  其中 $u=r^2$。所以
  \[
    \boxed{0<p<1\text{ 时收敛， }p\ge1\text{ 时发散}}.
  \]
\end{enumerate}
\end{xsolution}

\paragraph{第5题}
\begin{xsolution}
以下设 $p>0$。当 $0<p<1$ 时，一维奇积分 $\int\abs{y-c}^{-p}\,\dd y$ 在 $c\in\RR$ 时局部可积，而且对 $c$ 在有界集内可作一致控制。因此无论连续函数 $f$ 的图形如何与矩形 $D$ 相交，原二重广义积分均收敛。

当 $p\ge1$ 时，若存在 $x_0\in[a,A]$ 使
\[
 b<f(x_0)<B,
\]
则由连续性，某个 $x$ 的小区间内仍有 $f(x)\in(b,B)$；对这些 $x$，沿 $y=f(x)$ 存在一维不可积奇性，因此二重积分发散。

若 $f([a,A])$ 与开区间 $(b,B)$ 不相交，则由区间连通性，必有 $f\le b$ 或 $f\ge B$。若与边界还有正距离，积分显然收敛。只需讨论接触边界的情形。以 $f\le b$ 为例，令
\[
 h(x)=b-f(x)\ge0.
\]
把 $y=b+s$ 后，靠近奇点的主要部分为
\[
 \int_a^A\int_0^\delta(s+h(x))^{-p}\,\dd s\,\dd x.
\]
因此：
\begin{itemize}
  \item $p=1$ 时，收敛当且仅当 $\displaystyle\int_a^A\abs{\ln h(x)}\,\dd x<+\infty$（只需在 $h=0$ 的邻域检查）；
  \item $p>1$ 时，收敛当且仅当 $\displaystyle\int_a^Ah(x)^{1-p}\,\dd x<+\infty$。
\end{itemize}
$f\ge B$ 时令 $h(x)=f(x)-B$，结论完全相同。

所以一般连续 $f$ 下，$p\ge1$ 的答案确实依赖于 $f$ 接触矩形边界的方式；仅由题中条件不能进一步化成单一的 $p$ 范围。
\end{xsolution}

\paragraph{第6题}
\begin{xsolution}
\begin{enumerate}[label=(\arabic*)]
  \item 极坐标下
  \begin{align*}
   I&=2\pi\int_0^1r\ln\frac1r\,\dd r
   =2\pi\cdot\frac14
   =\boxed{\frac\pi2}.
  \end{align*}

  \item 区域为 $0\le x\le\pi,0\le y\le x$。令 $u=x-y$，先交换积分次序：
  \begin{align*}
   I&=\int_0^\pi\dd x\int_0^x\ln(\sin(x-y))\,\dd y\\
   &=\int_0^\pi(\pi-u)\ln(\sin u)\,\dd u.
  \end{align*}
  由 $\ln(\sin u)$ 关于 $\pi/2$ 对称，
  \[
   \int_0^\pi u\ln(\sin u)\,\dd u
   =\frac\pi2\int_0^\pi\ln(\sin u)\,\dd u.
  \]
  再用经典积分 $\int_0^\pi\ln(\sin u)\,\dd u=-\pi\ln2$，得到
  \[
    \boxed{I=-\frac{\pi^2}{2}\ln2}.
  \]
\end{enumerate}
\end{xsolution}

\subsection*{22.5.4 练习题}

\paragraph{第1题}
\begin{xsolution}
\begin{enumerate}[label=(\arabic*)]
  \item 记
  \[
   \Delta=\det
   \begin{pmatrix}
    a_1&b_1&c_1\\
    a_2&b_2&c_2\\
    a_3&b_3&c_3
   \end{pmatrix}\ne0.
  \]
  作线性变换
  \[
   u_i=a_ix+b_iy+c_iz,
   \qquad i=1,2,3.
  \]
  原立体变为长方体 $\abs{u_i}\le h_i$，并且
  \[
   \dd x\,\dd y\,\dd z=\frac1{\abs\Delta}\,\dd u_1\,\dd u_2\,\dd u_3.
  \]
  故体积
  \[
   \boxed{V=\frac{8h_1h_2h_3}{\abs\Delta}}.
  \]

  \item 用球坐标 $z=r\cos\varphi$。曲面方程化为
  \[
   r^4=a^3r\cos\varphi,
   \qquad r^3=a^3\cos\varphi.
  \]
  因而 $0\le\varphi\le\pi/2$，$0\le r\le a(\cos\varphi)^{1/3}$。于是
  \begin{align*}
   V&=\int_0^{2\pi}\int_0^{\pi/2}\int_0^{a(\cos\varphi)^{1/3}}
   r^2\sin\varphi\,\dd r\,\dd\varphi\,\dd\theta\\
   &=\frac{2\pi a^3}{3}\int_0^{\pi/2}\cos\varphi\sin\varphi\,\dd\varphi
   =\boxed{\frac{\pi a^3}{3}}.
  \end{align*}

  \item 令
  \[
   X=\frac xa,
   \qquad Y=\frac yb,
   \qquad Z=\frac zc.
  \]
  则体积乘上因子 $abc$，而边界变为
  \[
   (X^2+Y^2+Z^2)^2=X^2+Y^2.
  \]
  在球坐标中为 $r^2=\sin^2\varphi$，即 $0\le r\le\sin\varphi$。故
  \begin{align*}
   V&=abc\int_0^{2\pi}\int_0^\pi\int_0^{\sin\varphi}
   r^2\sin\varphi\,\dd r\,\dd\varphi\,\dd\theta\\
   &=\frac{2\pi abc}{3}\int_0^\pi\sin^4\varphi\,\dd\varphi
   =\boxed{\frac{\pi^2abc}{4}}.
  \end{align*}

  \item 采用柱坐标。由
  \[
   r^4+z^4=z
  \]
  知 $0\le z\le1$，且
  \[
   0\le r\le(z-z^4)^{1/4}.
  \]
  因而
  \begin{align*}
   V&=\pi\int_0^1\sqrt{z-z^4}\,\dd z
   =\pi\int_0^1z^{1/2}(1-z^3)^{1/2}\,\dd z\\
   &=\frac\pi3\int_0^1u^{-1/2}(1-u)^{1/2}\,\dd u
   \qquad(u=z^3)\\
   &=\frac{2\pi}{3}\int_0^{\pi/2}\cos^2 t\,\dd t
   \qquad(u=\sin^2t)\\
   &=\boxed{\frac{\pi^2}{6}}.
  \end{align*}
\end{enumerate}
\end{xsolution}

\paragraph{第2题}
\begin{xsolution}
\begin{enumerate}[label=(\arabic*)]
  \item 在球坐标中，曲面写成
  \[
   r=R(\varphi,\theta)=\sin\varphi\sqrt{\cos2\theta},
  \]
  其中 $0\le\varphi\le\pi$，而 $\cos2\theta\ge0$，即
  \[
   \theta\in\left[-\frac\pi4,\frac\pi4\right]
   \cup\left[\frac{3\pi}{4},\frac{5\pi}{4}\right].
  \]
  对球坐标中的径向图 $r=R(\varphi,\theta)$，
  \[
   \dd S=R^2\sin\varphi
   \sqrt{1+\left(\frac{R_\varphi}{R}\right)^2+
   \left(\frac{R_\theta}{R\sin\varphi}\right)^2}
   \,\dd\varphi\,\dd\theta.
  \]
  此处
  \[
   \frac{R_\varphi}{R}=\cot\varphi,
   \qquad
   \frac{R_\theta}{R}=-\tan2\theta,
  \]
  所以在 $\cos2\theta>0$ 的区域内恰有
  \[
   \dd S=\sin^2\varphi\,\dd\varphi\,\dd\theta.
  \]
  两个 $\theta$ 区间的总长度为 $\pi$，故
  \[
   \boxed{S=\pi\int_0^\pi\sin^2\varphi\,\dd\varphi=\frac{\pi^2}{2}}.
  \]

  \item 球坐标下方程为
  \[
   r^4=(r\cos\varphi)^3,
   \qquad r=\cos^3\varphi,
  \]
  因而 $0\le\varphi\le\pi/2$。径向图面积元为
  \[
   \dd S=\cos^5\varphi\sin\varphi
   \sqrt{\cos^2\varphi+9\sin^2\varphi}
   \,\dd\varphi\,\dd\theta.
  \]
  令 $s=\cos\varphi$，得
  \begin{align*}
   S&=2\pi\int_0^1s^5\sqrt{9-8s^2}\,\dd s
   =\boxed{\frac{313\pi}{560}}.
  \end{align*}

  \item 若母线以弧长 $s$ 为参数，则旋转面面积微元为 $2\pi y\,\dd s$，所以
  \[
   \boxed{S=2\pi\int_\Gamma y\,\dd s}.
  \]
  若进一步 $f\in C^1[a,b]$，则 $\dd s=\sqrt{1+[f'(x)]^2}\,\dd x$，从而
  \[
   \boxed{S=2\pi\int_a^bf(x)\sqrt{1+[f'(x)]^2}\,\dd x}.
  \]
\end{enumerate}
\end{xsolution}

\paragraph{第3题}
\begin{xsolution}
抛物面写成
\[
 z=\frac12r^2,
 \qquad0\le r\le\sqrt2.
\]
其面积元为
\[
 \dd S=\sqrt{1+z_x^2+z_y^2}\,\dd x\,\dd y
 =\sqrt{1+r^2}\,r\,\dd r\,\dd\theta.
\]
面密度 $\rho=z=r^2/2$，故质量
\begin{align*}
 m&=\int_0^{2\pi}\int_0^{\sqrt2}
 \frac{r^2}{2}\sqrt{1+r^2}\,r\,\dd r\,\dd\theta\\
 &=\pi\int_0^{\sqrt2}r^3\sqrt{1+r^2}\,\dd r
 =\boxed{\frac{2\pi}{15}(6\sqrt3+1)}.
\end{align*}
\end{xsolution}

\paragraph{第4题}
\begin{xsolution}
取圆盘中心为原点，细棒沿 $z$ 轴，近端为 $z=a$，远端为 $z=a+l$。圆盘面密度为 $\mu$。位于轴上距离圆盘为 $z>0$ 的单位质量所受引力大小为
\begin{align*}
 g(z)&=\int_0^R\frac{G\mu z}{(r^2+z^2)^{3/2}}\,2\pi r\,\dd r\\
 &=2\pi G\mu\left(1-\frac{z}{\sqrt{z^2+R^2}}\right).
\end{align*}
细棒线密度为 $\rho$，所以圆盘对整根细棒的引力大小为
\begin{align*}
 F&=\rho\int_a^{a+l}g(z)\,\dd z\\
 &=2\pi G\mu\rho
 \left[z-\sqrt{z^2+R^2}\right]_{a}^{a+l}\\
 &=\boxed{2\pi G\mu\rho
 \left(l+\sqrt{a^2+R^2}-\sqrt{(a+l)^2+R^2}\right)}.
\end{align*}
方向沿细棒指向圆盘。
\end{xsolution}

\paragraph{第5题}
\begin{xsolution}
取大圆盘中心为原点，小圆盘中心在 $(a/2,0)$；若小圆盘位于其他方向，最后把坐标轴旋转即可。原大圆盘的密度为到原点的距离 $r$，故总质量
\[
 M_0=\int_0^{2\pi}\int_0^a r\cdot r\,\dd r\,\dd\theta
 =\frac{2\pi a^3}{3},
\]
且重心在原点。

被挖去的小圆盘在原点极坐标中为
\[
 -\frac\pi2\le\theta\le\frac\pi2,
 \qquad0\le r\le a\cos\theta.
\]
其质量为
\[
 M_1=\int_{-\pi/2}^{\pi/2}\int_0^{a\cos\theta}r^2\,\dd r\,\dd\theta
 =\frac{4a^3}{9}.
\]
它关于 $y$ 轴的一阶矩为
\[
 Q_{y,1}=\int_{-\pi/2}^{\pi/2}\int_0^{a\cos\theta}
 (r\cos\theta)r^2\,\dd r\,\dd\theta
 =\frac{4a^4}{15},
\]
而关于 $x$ 轴的一阶矩为零。剩余部分的质量及一阶矩因此为
\[
 M=M_0-M_1=\frac{2a^3}{9}(3\pi-2),
 \qquad Q_y=-\frac{4a^4}{15},
 \qquad Q_x=0.
\]
所以重心坐标为
\[
 \boxed{
 \left(\bar x,\bar y\right)
 =\left(-\frac{6a}{5(3\pi-2)},0\right)}.
\]
若小圆盘中心沿单位向量 $\vect e$，则重心向量为
$-\dfrac{6a}{5(3\pi-2)}\vect e$。
\end{xsolution}

\paragraph{第6题}
\begin{xsolution}
设凸形物体为紧凸集 $K\subset\RR^3$，连续密度 $\rho\ge0$，总质量
\[
 m=\iiint_K\rho\,\dd V>0,
\]
重心为
\[
 C=\frac1m\iiint_K\vect x\,\rho(\vect x)\,\dd V.
\]
若 $C\notin K$，由凸集分离定理，存在向量 $\vect n$ 与常数 $\alpha$，使
\[
 \vect n\cdot C>\alpha\ge\vect n\cdot\vect x,
 \qquad\forall\vect x\in K.
\]
但由重心定义
\[
 \vect n\cdot C
 =\frac1m\iiint_K(\vect n\cdot\vect x)\rho(\vect x)\,\dd V
 \le\alpha,
\]
矛盾。因此 $C\in K$。

事实上，由 $m>0$ 与 $\rho\ge0$ 的连续性，存在物体内的一个相对开集使 $\rho>0$。若 $C$ 位于边界支撑平面 $\vect n\cdot\vect x=\alpha$ 上，则物体内有 $\vect n\cdot\vect x\le\alpha$，并在上述正密度开集的一部分上严格小于 $\alpha$，从而
\[
  \vect n\cdot C
  =\frac1m\iiint_K(\vect n\cdot\vect x)\rho(\vect x)\,\dd V<\alpha,
\]
与 $C$ 在该支撑平面上矛盾。因此 $C$ 不在边界上，故重心位于凸形物体内部。
\end{xsolution}

\paragraph{第7题}
\begin{xsolution}
对因子个数 $m$ 作归纳。$m=1$ 时由 $1/p_1=1$ 得 $p_1=1$，结论即为等式；$m=2$ 时就是 Hölder 不等式。

设结论对 $m-1$ 个因子成立，现考虑 $m\ge3$。令
\[
 \frac1r=\frac1{p_{m-1}}+\frac1{p_m}.
\]
由于
\[
 \sum_{i=1}^{m-2}\frac1{p_i}+\frac1r=1,
\]
有 $r>1$。又
\[
 \frac{r}{p_{m-1}}+\frac{r}{p_m}=1,
\]
故对 $u_{m-1},u_m$ 应用 Hölder 不等式可得
\begin{align*}
 \norm{u_{m-1}u_m}_r^r
 &=\iint_\Omega\abs{u_{m-1}}^r\abs{u_m}^r\,\dd x\,\dd y\\
 &\le
 \left(\iint_\Omega\abs{u_{m-1}}^{p_{m-1}}\,\dd x\,\dd y\right)^{r/p_{m-1}}
 \left(\iint_\Omega\abs{u_m}^{p_m}\,\dd x\,\dd y\right)^{r/p_m}.
\end{align*}
因此
\[
 \norm{u_{m-1}u_m}_r
 \le\norm{u_{m-1}}_{p_{m-1}}\norm{u_m}_{p_m}.
\]
再对 $m-1$ 个因子
\[
 u_1,\ldots,u_{m-2},u_{m-1}u_m
\]
以及指数 $p_1,\ldots,p_{m-2},r$ 使用归纳假设，得到
\begin{align*}
 \abs*{\iint_\Omega u_1u_2\cdots u_m\,\dd x\,\dd y}
 &\le\norm{u_1}_{p_1}\cdots\norm{u_{m-2}}_{p_{m-2}}
 \norm{u_{m-1}u_m}_r\\
 &\le\norm{u_1}_{p_1}\norm{u_2}_{p_2}\cdots\norm{u_m}_{p_m}.
\end{align*}
因此题中不带绝对值的积分当然也满足所给上界。
\end{xsolution}

\paragraph{第8题}
\begin{xsolution}
令
\[
 F(x)=\int_c^df(x,y)\,\dd y.
\]
则由 Fubini 交换积分次序以及关于 $x$ 的 Cauchy--Schwarz 不等式，
\begin{align*}
 \int_a^bF^2(x)\,\dd x
 &=\int_c^d\int_c^d
 \left(\int_a^bf(x,y)f(x,z)\,\dd x\right)\dd y\,\dd z\\
 &\le\int_c^d\int_c^d
 \left(\int_a^bf^2(x,y)\,\dd x\right)^{1/2}
 \left(\int_a^bf^2(x,z)\,\dd x\right)^{1/2}
 \dd y\,\dd z\\
 &=\left[
 \int_c^d\left(\int_a^bf^2(x,y)\,\dd x\right)^{1/2}\dd y
 \right]^2.
\end{align*}
两边开平方即得所要求的不等式。
\end{xsolution}

\subsection*{22.6.2 第一组参考题}

\paragraph{第1题}
\begin{xsolution}
先考察可积性。若 $(x_0,y_0)$ 的两个坐标都是无理数，则对任意 $N\in\NN$，可以取其充分小的邻域，使其中所有分母不超过 $N$ 的有理数都避开 $x_0,y_0$。于是该邻域内
\[
  0\le f(x,y)\le\frac1N,
\]
故 $f$ 在 $(x_0,y_0)$ 连续且值为 $0$。

因此 $f$ 的不连续点只能落在
\[
  \bigl((\QQ\cap[0,1])\times[0,1]\bigr)
  \cup
  \bigl([0,1]\times(\QQ\cap[0,1])\bigr),
\]
这是可列条直线之并，是零测度集。函数又有界，故由可积判别准则，$f$ 在 $D$ 上可积。并且连续点处函数值均为 $0$，由例题 22.1.2 的结论
\[
  \iint_Df(x,y)\,\dd x\,\dd y=0.
\]

下面说明两个二次积分均不存在。固定一个有理数 $y=p_y/q_y$。当 $x$ 为无理数时，
\[
  f(x,y)=\frac1{q_y},
\]
而当 $x$ 为有理数时 $f(x,y)=0$。因此 $x\mapsto f(x,y)$ 在 $[0,1]$ 上处处不连续，故
\[
  \int_0^1f(x,y)\,\dd x
\]
不存在。于是先对 $x$ 积分的二次积分不存在。同理，固定任一有理数 $x$，$y\mapsto f(x,y)$ 也处处不连续，所以先对 $y$ 积分的二次积分也不存在。
\end{xsolution}

\paragraph{第2题}
\begin{xsolution}
先考察截面。固定 $y\in[0,1]$。若 $y$ 是无理数或 $y=0$，则除至多有限情形外截面函数恒为 $0$；若 $y=p_y/q_y\ne0$ 为既约有理数，则
\[
  f(x,y)=1
\]
只可能发生在 $x$ 是分母同为 $q_y$ 的既约有理数时。这样的 $x\in[0,1]$ 只有有限个。因此对每个固定的 $y$，$x\mapsto f(x,y)$ 都只在有限点取 $1$，其余处取 $0$，故可积且积分为 $0$。于是
\[
  \int_0^1\dd y\int_0^1f(x,y)\,\dd x=0.
\]
完全同理，另一二次积分也存在且为 $0$。

但 $f$ 在 $D$ 的每一点都不连续。事实上，$f=0$ 的点在 $D$ 中稠密，例如至少有一个坐标为无理数的点构成稠密集。另一方面，$f=1$ 的点也稠密：给定任意 $(x_0,y_0)$ 和任意邻域，取充分大的素数 $q$，可以分别选取整数 $p_1,p_2$，使 $p_i/q$ 落在对应坐标的小邻域内；由于 $q$ 为素数并可避开端点，两个分数可取为既约分数且分母同为 $q$，故此点处 $f=1$。于是任一点的任意邻域内同时有函数值 $0$ 和 $1$，所以 $f$ 处处不连续。其不连续点集为整个正方形，非零测度，因此 $f$ 在 $D$ 上不可积。
\end{xsolution}

\paragraph{第3题}
\begin{xsolution}
\begin{enumerate}[label=(\arabic*)]
  \item 固定 $x$。当 $0\le x\le1/4$ 时，$xy\le1/4$ 对所有 $0\le y\le1$ 成立；当 $1/4<x\le1$ 时，变号点为 $y=1/(4x)$。故
  \begin{align*}
    A&=\int_0^{1/4}\int_0^1\left(\frac14-xy\right)\dd y\,\dd x\\
     &\quad+\int_{1/4}^1\left[
      \int_0^{1/(4x)}\left(\frac14-xy\right)\dd y
      +\int_{1/(4x)}^1\left(xy-\frac14\right)\dd y
    \right]\dd x\\
    &=\boxed{\frac3{32}+\frac18\ln2}.
  \end{align*}

  \item 由题设两式相减，
  \[
    \iint_D\left(xy-\frac14\right)f(x,y)\,\dd x\,\dd y=1.
  \]
  设
  \[
    M=\max_{(x,y)\in D}\abs{f(x,y)}.
  \]
  则
  \[
    1\le M\iint_D\abs*{xy-\frac14}\,\dd x\,\dd y=MA.
  \]
  因而 $M\ge1/A$。由于连续函数在紧集 $D$ 上达到最大绝对值，存在 $(\xi,\eta)\in D$ 使
  \[
    \boxed{\abs{f(\xi,\eta)}\ge\frac1A}.
  \]
\end{enumerate}
\end{xsolution}

\paragraph{第4题}
\begin{xsolution}
区域 $\abs{x}+\abs{y}\le10$ 是两条对角线均为 $20$ 的菱形，面积为
\[
  \abs{D}=200.
\]
又
\[
  100<100+\cos^2x+\cos^2y\le102
\]
除去零面积的特殊点外均严格成立相应的一侧，因此
\[
  \frac{200}{102}<
  \iint_D\frac{\dd x\,\dd y}{100+\cos^2x+\cos^2y}
  <\frac{200}{100}.
\]
由于
\[
  \frac{200}{102}=\frac{100}{51}>\frac{49}{25}=1.96,
\]
所以
\[
  \boxed{1.96<\iint_D\frac{\dd x\,\dd y}{100+\cos^2x+\cos^2y}<2}.
\]
\end{xsolution}

\paragraph{第5题}
\begin{xsolution}
令
\[
  k=\sqrt{a^2+b^2+c^2}.
\]
当 $k=0$ 时结论直接成立。以下设 $k>0$。取正交变换，把向量 $(a,b,c)$ 变为新坐标系中第三坐标轴方向的向量 $(0,0,k)$。正交变换保持单位球和体积元不变，因此原积分等于
\[
  \iiint_{u^2+v^2+w^2\le1}f(kw)\,\dd u\,\dd v\,\dd w.
\]
固定 $w\in[-1,1]$，截面 $u^2+v^2\le1-w^2$ 的面积为 $\pi(1-w^2)$。于是
\[
  \boxed{
  \iiint_{x^2+y^2+z^2\le1}f(ax+by+cz)\,\dd x\,\dd y\,\dd z
  =\pi\int_{-1}^1(1-u^2)f(ku)\,\dd u }.
\]
\end{xsolution}

\paragraph{第6题}
\begin{xsolution}
记
\[
  L=\max_{x\in[a,b]}\bigl(\psi(x)-\varphi(x)\bigr)<+\infty.
\]
对固定 $x$ 及 $y\in[\varphi(x),\psi(x)]$，由边界条件 $f(x,\varphi(x))=0$，
\[
  f(x,y)=\int_{\varphi(x)}^y f_y(x,t)\,\dd t.
\]
由 Cauchy--Schwarz 不等式
\[
  \abs{f(x,y)}^2
  \le(y-\varphi(x))
  \int_{\varphi(x)}^y f_y^2(x,t)\,\dd t
  \le L\int_{\varphi(x)}^{\psi(x)}f_y^2(x,t)\,\dd t.
\]
再对 $y$ 从 $\varphi(x)$ 到 $\psi(x)$ 积分，得到
\[
  \int_{\varphi(x)}^{\psi(x)}f^2(x,y)\,\dd y
  \le L^2\int_{\varphi(x)}^{\psi(x)}f_y^2(x,y)\,\dd y.
\]
最后对 $x$ 积分，
\[
  \iint_D f^2(x,y)\,\dd x\,\dd y
  \le L^2\iint_D f_y^2(x,y)\,\dd x\,\dd y.
\]
故可取
\[
  \boxed{K=L^2}.
\]
\end{xsolution}

\paragraph{第7题}
\begin{xsolution}
先证 Minkowski 不等式。$p=1$ 时由点态三角不等式积分即得。以下设 $p>1$，令 $q=p/(p-1)$。由
\[
  \abs{u+v}^p\le\abs{u}\,\abs{u+v}^{p-1}
  +\abs{v}\,\abs{u+v}^{p-1}
\]
及 Hölder 不等式，
\begin{align*}
  \norm{u+v}_p^p
  &\le \norm{u}_p
  \left(\iint_\Omega\abs{u+v}^{(p-1)q}\right)^{1/q}\\
  &\quad+\norm{v}_p
  \left(\iint_\Omega\abs{u+v}^{(p-1)q}\right)^{1/q}\\
  &=(\norm{u}_p+\norm{v}_p)\norm{u+v}_p^{p-1}.
\end{align*}
若 $\norm{u+v}_p>0$，约去最后一个因子即得
\[
  \boxed{\norm{u+v}_p\le\norm{u}_p+\norm{v}_p}.
\]
零范数情形显然成立。

Hölder 不等式在非零函数情形取等号的条件是：存在常数 $c>0$，使
\[
  \abs{u}^p=c\abs{v}^q
\]
几乎处处成立；若讨论不带绝对值的 $\iint uv$，还需 $uv$ 几乎处处同号。

Minkowski 不等式在 $p>1$ 时取等号的条件是 $u,v$ 几乎处处为同向比例，即存在 $\lambda\ge0$ 使
\[
  u=\lambda v
\]
几乎处处成立（包含其中一个函数为零的退化情形）。当 $p=1$ 时，取等号当且仅当
\[
  u(x,y)v(x,y)\ge0
\]
几乎处处成立。
\end{xsolution}

\paragraph{第8题}
\begin{xsolution}
令
\[
  w(x)=\frac12x(1-x).
\]
则 $w(0)=w(1)=0$，$w''(x)=-1$，并且
\[
  \int_0^1w(x)\,\dd x=\frac1{12}.
\]
固定 $y$。由于 $f(0,y)=f(1,y)=0$，两次分部积分给出
\[
  \int_0^1f(x,y)\,\dd x
  =-\int_0^1w(x)f_{xx}(x,y)\,\dd x.
\]
由于 $f(x,0)=f(x,1)=0$ 对一切 $x$ 成立，对 $x$ 求两次导数可得
\[
  f_{xx}(x,0)=f_{xx}(x,1)=0.
\]
再关于 $y$ 使用同一公式，
\[
  \iint_Df(x,y)\,\dd x\,\dd y
  =\iint_Dw(x)w(y)f_{xxyy}(x,y)\,\dd x\,\dd y.
\]
于是
\begin{align*}
 \abs*{\iint_Df(x,y)\,\dd x\,\dd y}
 &\le B\int_0^1w(x)\,\dd x\int_0^1w(y)\,\dd y\\
 &=\boxed{\frac{B}{144}}.
\end{align*}
\end{xsolution}

\paragraph{第9题}
\begin{xsolution}
设 $f$ 在 $D$ 上可积且
\[
  I=\iint_Df(x,y)\,\dd x\,\dd y>0.
\]
取足够细的分划 $T$，使对应的下和 $L(f,T)>0$。于是至少存在一个闭子区域 $U$，其面积为正，且该子区域上的下确界 $m_U$ 满足
\[
  m_U>0;
\]
否则所有小区域的下确界均不大于 $0$，下和不可能为正。故对所有 $(x,y)\in U$，
\[
  f(x,y)\ge m_U>0.
\]
所以存在所要求的闭子区域 $U$。
\end{xsolution}

\paragraph{第10题}
\begin{xsolution}
设 $\abs\Omega$ 表示 $\Omega$ 的体积，且 $\abs\Omega>0$。连续函数在紧集 $\Omega$ 上达到最小值与最大值，记为 $m,M$。于是
\[
  m\abs\Omega\le\int\cdots\int_\Omega f\le M\abs\Omega.
\]
因此积分平均值
\[
  A=\frac1{\abs\Omega}\int\cdots\int_\Omega f
\]
满足 $m\le A\le M$。区域 $\Omega$ 是连通的，故连续像 $f(\Omega)$ 是区间，并包含 $[m,M]$ 中的每个中间值。因此存在 $\xi\in\Omega$ 使 $f(\xi)=A$，即
\[
  \boxed{\int\cdots\int_\Omega f=f(\xi)\abs\Omega}.
\]
\end{xsolution}

\paragraph{第11题}
\begin{xsolution}
记
\[
  P=\int_a^bp(x)\,\dd x,
  \quad F=\int_a^bp(x)f(x)\,\dd x,
  \quad G=\int_a^bp(x)g(x)\,\dd x.
\]
则
\begin{align*}
 &P\int_a^bp(x)f(x)g(x)\,\dd x-FG\\
 &\quad=\frac12\int_a^b\int_a^b
 p(x)p(y)[f(x)-f(y)][g(x)-g(y)]\,\dd x\,\dd y.
\end{align*}
由于 $f,g$ 都单调增加，括号中的两个差总是同号，且 $p\ge0$，所以右端非负。于是
\[
  \left(\int_a^bp f\right)\left(\int_a^bp g\right)
  \le\left(\int_a^bp\right)\left(\int_a^bpfg\right).
\]
\end{xsolution}

\paragraph{第12题}
\begin{xsolution}
使用上一题的双重积分恒等式，但令权函数为正函数 $p(x)=f(x)$，并取
\[
  F_1(x)=x,
  \qquad G_1(x)=f(x).
\]
$F_1$ 单调增加而 $G_1$ 单调减少，因此
\[
 [F_1(x)-F_1(y)][G_1(x)-G_1(y)]\le0.
\]
于是相应的 Chebyshev 不等式反向：
\[
 \left(\int_0^1f\right)
 \left(\int_0^1x f^2\right)
 \le
 \left(\int_0^1x f\right)
 \left(\int_0^1 f^2\right).
\]
因 $f>0$，两个分母均为正，约去后即得
\[
  \boxed{
  \frac{\int_0^1x f^2(x)\,\dd x}{\int_0^1x f(x)\,\dd x}
  \le
  \frac{\int_0^1f^2(x)\,\dd x}{\int_0^1f(x)\,\dd x}}.
\]
\end{xsolution}

\paragraph{第13题}
\begin{xsolution}
记球心为 $O$，固定点 $P=(a,b,c)$，则 $OP=A>R$。对球内任意点 $X$，三角不等式给出
\[
  A-R\le PX\le A+R.
\]
因此
\[
  \frac1{A+R}\le\frac1{PX}\le\frac1{A-R}.
\]
在半径 $R$ 的球上积分，并使用球体积 $4\pi R^3/3$，得到
\[
  \boxed{
  \frac{4\pi}{3}\frac{R^3}{A+R}
  \le I\le
  \frac{4\pi}{3}\frac{R^3}{A-R}}.
\]
\end{xsolution}

\paragraph{第14题}
\begin{xsolution}
由变上限三重积分求导，三个面的贡献相同，故
\[
  F'(t)=3\int_0^t\int_0^t f(tyz)\,\dd y\,\dd z.
\]
设
\[
  J(t)=\int_0^t\int_0^t f(tyz)\,\dd y\,\dd z.
\]
固定 $y$ 并令 $s=tyz$，则
\[
  \int_0^tf(tyz)\,\dd z
  =\frac{g(t^2y)}{ty},
\]
其中 $g(u)=\int_0^uf(s)\,\dd s$，在 $y=0$ 处按连续极限理解。因此
\[
  J(t)=\frac1t\int_0^t\frac{g(t^2y)}y\,\dd y.
\]
令 $u=t^2y$，得到
\[
  J(t)=\frac1t\int_0^{t^3}\frac{g(u)}u\,\dd u.
\]
故
\[
  \boxed{F'(t)=\frac3t\int_0^{t^3}\frac{g(u)}u\,\dd u}.
\]
\end{xsolution}

\paragraph{第15题}
\begin{xsolution}
由于 $s$ 是逆时针弧长参数，单位切向量为
\[
  (f'(s),\varphi'(s))=(\cos\theta(s),\sin\theta(s)).
\]
对于逆时针定向的凸椭圆，外单位法向量为
\[
  (\sin\theta(s),-\cos\theta(s)).
\]
因此：
\begin{enumerate}[label=(\arabic*)]
  \item 距离曲线 $\Gamma$ 为 $t$ 的外侧点可写成
  \[
    \boxed{
    x=f(s)+t\sin\theta(s),
    \qquad
    y=\varphi(s)-t\cos\theta(s)},
  \]
  其中 $0\le s\le2\pi l$，$0<t<l$。

  \item 计算 Jacobi 行列式。利用 $f'=\cos\theta$、$\varphi'=\sin\theta$，
  \begin{align*}
    \frac{\partial(x,y)}{\partial(s,t)}
    &=\det
    \begin{pmatrix}
      \cos\theta+t\theta'\cos\theta&\sin\theta\\
      \sin\theta+t\theta'\sin\theta&-\cos\theta
    \end{pmatrix}\\
    &=-(1+t\theta'(s)).
  \end{align*}
  椭圆严格凸，所以 $\theta'(s)>0$，故面积元为
  \[
    \dd x\,\dd y=(1+t\theta'(s))\,\dd s\,\dd t.
  \]
  又切向量绕一周，$\theta$ 总增加 $2\pi$，于是
  \begin{align*}
    \abs{D}
    &=\int_0^{2\pi l}\int_0^l(1+t\theta'(s))\,\dd t\,\dd s\\
    &=2\pi l^2+\frac{l^2}{2}\int_0^{2\pi l}\theta'(s)\,\dd s\\
    &=2\pi l^2+\pi l^2
    =\boxed{3\pi l^2}.
  \end{align*}
\end{enumerate}
\end{xsolution}

\subsection*{22.6.2 第二组参考题}

\paragraph{第1题}
\begin{xsolution}
由 $1/p+1/q=1$ 且 $p,q>0$，实际有 $p,q>1$。对 Young 不等式
\[
  XY\le\frac{X^p}{p}+\frac{Y^q}{q}
\]
取
\[
  X=\varepsilon^{1/p}a,
  \qquad
  Y=\varepsilon^{-1/p}b.
\]
则 $XY=ab$，且
\[
  X^p=\varepsilon a^p,
  \qquad
  Y^q=\varepsilon^{-q/p}b^q.
\]
故
\[
  ab\le\frac{\varepsilon a^p}{p}
  +\frac{\varepsilon^{-q/p}b^q}{q}.
\]
又 $1/p<1,1/q<1$，两项均非负，因此
\[
  \frac{\varepsilon a^p}{p}
  +\frac{\varepsilon^{-q/p}b^q}{q}
  \le\varepsilon a^p+\varepsilon^{-q/p}b^q.
\]
\end{xsolution}

\paragraph{第2题}
\begin{xsolution}
由例题 22.5.9，存在 $0<\lambda<1$ 使
\[
  \norm{u}_q\le\norm{u}_p^\lambda\norm{u}_r^{1-\lambda},
  \qquad
  \frac1q=\frac\lambda p+\frac{1-\lambda}{r}.
\]
由此解得
\[
  \frac{1-\lambda}{\lambda}
  =\frac{\frac1p-\frac1q}{\frac1q-\frac1r}
  =\mu.
\]
在上一题中取共轭指数
\[
  \alpha=\frac1{1-\lambda},
  \qquad
  \beta=\frac1\lambda,
\]
并令
\[
  a=\norm{u}_r^{1-\lambda},
  \qquad
  b=\norm{u}_p^\lambda.
\]
则对任意 $\varepsilon>0$，上一题的较粗上界给出
\[
  \norm{u}_r^{1-\lambda}\norm{u}_p^\lambda
  \le\varepsilon\norm{u}_r
  +\varepsilon^{-\beta/\alpha}\norm{u}_p.
\]
而
\[
  \frac\beta\alpha=\frac{1-\lambda}{\lambda}=\mu.
\]
因此
\[
  \boxed{\norm{u}_q\le\varepsilon\norm{u}_r
  +\varepsilon^{-\mu}\norm{u}_p}.
\]
\end{xsolution}

\paragraph{第3题}
\begin{xsolution}
记 $A=\abs\Omega$。
\begin{enumerate}[label=(\arabic*)]
  \item 令
  \[
    M=\max_\Omega u.
  \]
  显然 $\Phi_p(u)\le M$。另一方面，对任意 $\delta>0$，由连续性，集合
  \[
    E_\delta=\{(x,y)\in\Omega\mid u(x,y)>M-\delta\}
  \]
  含有一个正面积邻域，故 $\abs{E_\delta}>0$。于是
  \[
    \Phi_p(u)
    \ge(M-\delta)
    \left(\frac{\abs{E_\delta}}A\right)^{1/p}.
  \]
  令 $p\to+\infty$ 得 $\liminf\Phi_p(u)\ge M-\delta$，再令 $\delta\to0$，结合上界可得
  \[
    \boxed{\lim_{p\to+\infty}\Phi_p(u)=M}.
  \]

  \item 因 $u$ 连续且恒正，$v=1/u$ 也连续恒正。令 $q=-p\to+\infty$，则
  \[
    \Phi_p(u)=\frac1{\Phi_q(v)}.
  \]
  由 (1)，
  \[
    \lim_{q\to+\infty}\Phi_q(v)=\max_\Omega\frac1u
    =\frac1{\min_\Omega u}.
  \]
  因此
  \[
    \boxed{\lim_{p\to-\infty}\Phi_p(u)=\min_\Omega u}.
  \]

  \item 记 $h=\ln u$。由于 $u$ 在紧集上有正的最小值，$h$ 有界。于是当 $p\to0$ 时，一致地有
  \[
    \ee^{ph}=1+ph+\bigO(p^2).
  \]
  因而
  \[
    \frac1A\iint_\Omega u^p
    =1+p\frac1A\iint_\Omega\ln u+\bigO(p^2).
  \]
  取对数并除以 $p$，
  \[
    \ln\Phi_p(u)
    =\frac1p\ln\left(\frac1A\iint_\Omega u^p\right)
    \longrightarrow\frac1A\iint_\Omega\ln u.
  \]
  指数化即得
  \[
    \boxed{\lim_{p\to0}\Phi_p(u)
    =\exp\left\{\frac1A\iint_\Omega\ln u\,\dd x\,\dd y\right\}}.
  \]
\end{enumerate}
\end{xsolution}

\paragraph{第4题}
\begin{xsolution}
设球心为 $O$，以向量 $\vect a=\overrightarrow{OP_0}$ 表示定点，记 $d=\norm{\vect a}<R$。球面点写成 $Q=R\vect n$，其中 $\norm{\vect n}=1$。$Q$ 处切平面为
\[
  \vect n\cdot\vect x=R.
\]
从 $P_0$ 向此平面作垂线，垂足为
\[
  P=P_0+(R-\vect a\cdot\vect n)\vect n.
\]
所以以 $P_0$ 为原点，所得封闭曲面的径向函数为
\[
  \rho(\vect n)=R-\vect a\cdot\vect n>0.
\]

\begin{enumerate}[label=(\arabic*)]
  \item 星形体的体积为
  \[
    V=\frac13\int_{S^2}\rho(\vect n)^3\,\dd\omega.
  \]
  展开并利用球面对称性
  \[
    \int_{S^2}\vect a\cdot\vect n\,\dd\omega=0,
    \qquad
    \int_{S^2}(\vect a\cdot\vect n)^2\,\dd\omega
    =\frac{4\pi}{3}d^2,
  \]
  得
  \begin{align*}
    V&=\frac13\left(4\pi R^3+3R\frac{4\pi}{3}d^2\right)\\
     &=\boxed{\frac{4\pi R}{3}(R^2+d^2)}.
  \end{align*}

  \item 把 $V$ 视为 $\vect a$ 的函数，
  \[
    \nabla_{\vect a}V=\frac{8\pi R}{3}\vect a.
  \]
  因而当 $P_0\ne O$ 时，在给定相同移动速度大小的条件下，体积增加率最大的方向是
  \[
    \boxed{\text{沿 }OP_0\text{ 从球心指向 }P_0\text{ 的方向}}.
  \]
  若 $P_0=O$，则一阶变化率在所有方向上均为 $0$。
\end{enumerate}
\end{xsolution}

\paragraph{第5题}
\begin{xsolution}
记
\[
  J(a)=\int_0^a\ee^{-x^2}\,\dd x.
\]
则
\[
  J(a)^2=\iint_{[0,a]^2}\ee^{-(x^2+y^2)}\,\dd x\,\dd y.
\]

半径为 $a$ 的第一象限四分之一圆盘包含在正方形内，所以
\begin{align*}
  J(a)^2
  &\ge\int_0^{\pi/2}\int_0^a\ee^{-r^2}r\,\dd r\,\dd\theta\\
  &=\frac\pi4(1-\ee^{-a^2}).
\end{align*}
这给出所需下界。

对上界，在极坐标中正方形沿方向 $\theta$ 的边界半径为
\[
  R(\theta)=\frac{a}{\max(\cos\theta,\sin\theta)}.
\]
因而
\[
  J(a)^2=\frac12\int_0^{\pi/2}
  \left(1-\ee^{-R(\theta)^2}\right)\dd\theta.
\]
函数 $\phi(s)=1-\ee^{-s}$ 在 $s\ge0$ 上为凹函数。又由正方形面积
\[
  a^2=\frac12\int_0^{\pi/2}R(\theta)^2\,\dd\theta,
\]
故
\[
  \frac2\pi\int_0^{\pi/2}R(\theta)^2\,\dd\theta
  =\frac{4a^2}{\pi}.
\]
由 Jensen 不等式
\[
  J(a)^2\le\frac\pi4
  \left(1-\ee^{-4a^2/\pi}\right).
\]
两边开平方，得到
\[
  \boxed{
  \frac{\sqrt\pi}{2}(1-\ee^{-a^2})^{1/2}
  \le J(a)\le
  \frac{\sqrt\pi}{2}(1-\ee^{-4a^2/\pi})^{1/2}}.
\]
\end{xsolution}

\paragraph{第6题}
\begin{xsolution}
把 $[v_1,v_2]$ 分成
\[
  v_1=t_0<t_1<\cdots<t_n=v_2.
\]
令
\[
  E_i=\{(x,y)\mid t_{i-1}\le f(x,y)\le t_i\}
\]
为相邻两条等位线之间的薄层。按照 $F$ 的定义，
\[
  \abs{E_i}=F(t_i)-F(t_{i-1}).
\]
在 $E_i$ 上有
\[
  t_{i-1}\le f(x,y)\le t_i.
\]
因而
\[
  t_{i-1}[F(t_i)-F(t_{i-1})]
  \le\iint_{E_i}f(x,y)\,\dd x\,\dd y
  \le t_i[F(t_i)-F(t_{i-1})].
\]
对 $i$ 求和。当分割的模趋于 $0$ 时，上、下两端均趋于 Riemann--Stieltjes 积分
\[
  \int_{v_1}^{v_2}v\,\dd F(v).
\]
题设又假定 $F$ 可微且 $F'$ 可积，因此
\[
  \int_{v_1}^{v_2}v\,\dd F(v)
  =\int_{v_1}^{v_2}vF'(v)\,\dd v.
\]
故
\[
  \boxed{
  \iint_{S(v_1,v_2)}f(x,y)\,\dd x\,\dd y
  =\int_{v_1}^{v_2}vF'(v)\,\dd v}.
\]
这就是 Catalan 公式。
\end{xsolution}
